\section{Methodology and experimental design}\label{sec:methods}

\subsection{Design principles}
Fibre degradation experiments are slow, and many things can change in a lab over several weeks. The design
therefore follows five rules:
\begin{enumerate}[label=\textbf{P\arabic*}]
  \item \textbf{Parallel, not sequential.} All test atmospheres run \emph{at the same time} from
        one light source, in identical optical arms. Laser drifts, lab humidity and seasonal
        effects then affect all arms equally and cancel in the comparison.
  \item \textbf{Change one factor at a time.} The decisive pair of arms differs \emph{only} in
        the presence of \Oxy{}: clean dry synthetic air versus clean dry \Ntwo{}, both passed through the
        same purifier and particle filter (\cref{tab:arms}).
  \item \textbf{Separate the surface from the alignment.} Throughput is recorded before and after
        re-peaking the input coupling, and the back-reflection from the input face is logged
        continuously. Alignment drift (H6) can then be told apart from changes of the end face.
  \item \textbf{Locate the loss.} At the end of every run a fixed sequence of non-destructive
        and destructive steps (\cref{sec:ladder}) splits the loss between input face, output face
        and fibre.
  \item \textbf{Decide in advance.} Metrics, the analysis and the decision thresholds
        (\cref{sec:analysis}) are written down before the confirmatory campaign starts. This
        guards against reading patterns into noisy data after the fact.
\end{enumerate}

\subsection{Phase 0: forensics of fibres already degraded in the lab}\label{sec:forensics}
The cheapest information is already in the lab. In the first two months the student collects
degraded fibres from the group's experiments (ideally 3--6, with as much history as possible:
power, hours of use, wavelength, environment, cleaning and polishing history, PC or APC ends) and:
\begin{enumerate}
  \item images both end faces \emph{before any cleaning} with the fibre inspection microscope (core
        zone, standard magnification and illumination, \cref{app:sops});
  \item measures throughput at low power (\SI{\le 0.3}{mW}) to avoid changing the surfaces further;
  \item runs the cleaning sequence of \cref{sec:ladder} on half of the fibres and keeps the others
        uncleaned for surface analysis (optical and fluorescence microscopy, SEM/EDX, AFM or
        white-light interferometry, Raman, and where available ToF-SIMS or XPS).
\end{enumerate}
In parallel, the student surveys likely contaminant sources near the group's fibre ends: silicone
grease, silicone tubing and adhesives, thermal paste, cable jackets, printed parts and cleaning
practices. Within weeks this shows whether the damage is a \emph{deposit} (and whether it is rich in carbon
or in silicon), a \emph{surface modification} or \emph{nothing visible}. That result already favours some hypotheses
over others and decides which test atmospheres are most informative.

\subsection{Test fibres and end-face preparation}
\begin{itemize}
  \item \textbf{Fibre type.} The same fibre as used in the experiments, so that the results transfer
        directly: Thorlabs P3-375PM-FC patch cables, FC/APC at both ends, built on PM-S350-HP
        pure-silica-core PANDA fibre (\cref{sec:sota-fibres}), ideally from a single batch. Thorlabs
        sells the bare fibre no longer, but a fibre with the same part number and specification is
        available from Coherent (Nufern). Cheaper non-PM SM300 cables, which also have a pure-silica
        core, can serve as extra arms where polarisation does not matter.
  \item \textbf{Format.} \SIrange{1}{2}{m} lengths, connectorised and polished in-house with one
        documented procedure and one epoxy (or bought as identical patch cables). The fibre is coiled
        loosely (bend radius well above the specified minimum) and taped to the table, so that bend
        loss cannot change.
  \item \textbf{Acceptance.} Every end face is inspected and imaged after polishing, and its geometry
        is measured if an end-face interferometer is available. Its initial coupling efficiency is measured
        at low power. Ends that miss the acceptance criteria (no defects in the core zone; coupling
        within a few per cent of the batch median) are re-polished before use.
  \item \textbf{Surface preparation is a variable.} At \SI{405}{nm}, cleaved faces degraded four times
        more slowly than glued-and-polished ones~\cite{Gonschior2013PhD}. Every test end is therefore
        prepared by exactly the same polishing procedure. A comparison of preparation methods (e.g.\ an
        additional final polishing step) is an optional arm in Campaign~II.
  \item \textbf{Traceability.} Each fibre gets an ID. Each end has its own polishing record and image
        set. The assignment of fibres to arms is randomised, and arm positions are swapped between
        runs so that any arm-specific optical problem cannot look like an atmosphere effect.
\end{itemize}

\subsection{Light source and multi-arm testbed}\label{sec:testbed}
\Cref{fig:testbed} shows the testbed. Its main elements are:

\paragraph{Source} A dedicated \SI{375}{nm} single-mode laser diode (\cref{sec:sota-hardware}) with
an optical isolator and anamorphic beam shaping. It frees the group's \SI{369.5}{nm} lasers for
experiments and allows long, uninterrupted runs. The photon energies at \SI{375}{nm} and
\SI{369.5}{nm} differ by only \SI{1.5}{\percent}, and no known silica absorption band lies between them.
Equivalent behaviour is therefore expected, but this assumption is \emph{tested} in Phase~4, which
runs at \SI{369.5}{nm}. If no diode is available, a pick-off of a few tens of milliwatts from an
existing \SI{369.5}{nm} system can feed the testbed.

\paragraph{Power distribution} A cascade of half-wave plates and polarising beam splitters divides the
beam into 4--6 arms with adjustable and equal powers. A further half-wave plate in each arm aligns the
polarisation to the slow axis of the PM fibre. The operational test power is \SI{3}{mW} guided; the
accelerated condition is \SI{10}{mW}. A sentinel arm at \SI{\le 0.1}{mW} in \Ntwo{} measures the
drift of the measurement system itself.

\paragraph{Coupling} An aspheric lens with near-UV anti-reflection coating ($f\approx\SIrange{4.5}{8}{mm}$,
chosen for mode matching, \cref{app:estimates}) focuses the beam through the cell window onto the
fibre end face, at the \ang{3.8} incidence an \ang{8} APC face requires. Commercial launch systems for
FC/APC (e.g.\ Thorlabs FiberPort PAF2A-5A or -7A) or a moulded asphere (e.g.\ C230TMD-A) can be used.
The high-index glasses of such aspheres transmit less well in the near UV, so their transmission at
\SI{369.5}{nm} should be checked. The lens, or the cell, sits on a stable XYZ mount. Automated peak-up with motorised mirror
mounts is a stretch goal; the baseline is a manual re-peak once a day with a fixed procedure.

\paragraph{Detection} A wedged pick-off in each arm feeds a reference photodiode ($P_\text{ref}$) and a
second photodiode that records the light reflected by the input face ($P_\text{back}$; with FC/APC
connectors this detector sits in the specularly reflected beam, which does not retrace the input). The output power $P_\text{out}$
is measured with an integrating-sphere or diffuser-fronted photodiode, which is insensitive to small
changes of the output beam. Silicon photodiodes lose UV responsivity and gain dark current under
prolonged UV exposure~\cite{HamamatsuS1336,Korde1987}. All detectors are therefore operated at low
intensity (attenuated, diffused or shuttered), photodiodes sold for UV resistance are preferred (e.g.\
Hamamatsu S12698 series~\cite{HamamatsuS12698}), and a reference power meter that is otherwise kept
in a drawer cross-checks them weekly. SiC photodiodes, the usual choice for UV-stable detection, are not
suitable here: \SI{369.5}{nm} lies at the edge of their spectral response.

\paragraph{Logging} All photodiodes, laser current and temperature, lab and arm temperatures, humidity,
pressure, flows and \Oxy{} readings are logged at \SI{1}{Hz} (stored as 1-minute means) by a Python
data logger (\cref{app:logging}). Every manual intervention is recorded in an event log with a
timestamp.

\subsection{Test atmospheres}\label{sec:atmospheres}
\begin{table}[t]
\centering\small
\caption{Test atmospheres. Core arms run in Campaign~I; extended arms are selected for Campaign~II
on the basis of Campaign~I and the Phase~0 forensics.}
\label{tab:arms}
\begin{tabularx}{\textwidth}{@{}l L{4.2cm} Y@{}}
\toprule
Arm & Atmosphere at both end faces & Purpose \\
\midrule
\multicolumn{3}{@{}l}{\emph{Core arms}}\\
\arm{AIR}  & ambient lab air, ends open as in the experiments & control; reproduces the problem \\
\arm{CDA}  & synthetic air (\SI{20.9}{\percent} \Oxy{} in \Ntwo{}, zero grade), purified, filtered, dry
           & contains \Oxy{} but no contaminants, particles or water \\
\arm{N2}   & \Ntwo{} (grade 5.0 or LN$_2$ boil-off), same purifier and filter, dry
           & identical to \arm{CDA} except for \Oxy{} $\Rightarrow$ decisive test of H1 vs H2 \\
\arm{ABS}  & sealed printed cell, pre-purged with \Ntwo{}, iron-based \Oxy{} absorber, no flow
           & the low-cost, passive solution to be deployed \\
\midrule
\multicolumn{3}{@{}l}{\emph{Extended arms}}\\
\arm{O2-x} & \Ntwo{} with $x=$ 10, 100, \num{1000}, \num{10000}\,ppm \Oxy{}
           & dose--response, sets the \Oxy{} specification for the cells \\
\arm{WET-N2}  & \Ntwo{} humidified to \SIrange{30}{50}{\percent} RH & water without \Oxy: assists damage (H3)
           or protects against carbon deposits (H2)~\cite{Weiller2010,Parthasarathy2024} \\
\arm{WET-CDA} & synthetic air at lab-like RH, filtered & lab air minus contaminants and particles (H2, H4) \\
\arm{FG}   & forming gas, \SI{5}{\percent} \Htwo{} in \Ntwo{} & hydrogen passivation of defects (H5) \\
\arm{N2+P} & \Ntwo{} cell containing a coupon of the printing material
           & are printed parts a contamination source? (H2) \\
\arm{+SIL} & a silicone coupon in one \arm{CDA} and one \arm{N2} cell
           & oxygen-assisted siloxane deposition (H1)~\cite{Fujifilmpatent2006} \\
\arm{VAC}  & vacuum (optional, if a chamber with fibre feedthrough is free) & link to in-vacuum fibre ends \\
\bottomrule
\end{tabularx}
\end{table}

\Cref{tab:arms} lists the atmospheres. Purged arms receive \SIrange{5}{20}{sccm} per cell through a
common chain of regulator, oxygen and moisture purifier and \SI{3}{nm} particle filter. This is far less
than the \SIrange{0.5}{1}{L/min} that Marciniak et al.\ blew through open collimators to keep particles
out~\cite{Marciniak2017}, but a sealed cell needs no more than a small positive outflow (\cref{app:estimates}). Short runs of
stainless-steel or nylon tubing are used, because polyurethane and fluoropolymer tubing let enough
oxygen through to matter at the ppm level (\cref{app:estimates}). Each outlet passes a humidity sensor, a
VOC sensor and (for \Oxy{}-free arms) an \Oxy{} sensor before reaching an oil bubbler or a long thin
exhaust line that stops back-diffusion. \Oxy{} mixtures for \arm{O2-x} are made by adding synthetic
air to \Ntwo{} with two flow controllers, or are bought as certified mixtures. Humidified arms use a
water bubbler in a split stream. The terms ``getter'' and ``absorber'' are used carefully. Vacuum
non-evaporable getters are not suitable at atmospheric pressure. The cells use iron-based oxygen
absorbers of the kind used in food and museum conservation (\cref{sec:sota-atmosphere}).

\subsection{End-face cells}\label{sec:cells}
Each fibre end that is not in ambient air sits in a small cell (\cref{fig:cell}). The requirements are:
\begin{enumerate}
  \item \textbf{Optical access} through a UV fused-silica window, AR-coated near \SI{370}{nm} and
        slightly wedged and tilted to avoid etalons and back-reflections into the fibre. The window lies
        in the converging beam well away from the focus, so the intensity on it is low. For a
        \SI{1}{mm} window at NA~0.1 the added spherical aberration is negligible
        (\cref{app:estimates}).
  \item \textbf{A small, swept volume} (\SIrange{10}{30}{cm^3}). Gas enters next to the end face and
        leaves downstream, so that anything released by the connector, epoxy or cable is carried
        \emph{away} from the face.
  \item \textbf{Tightness.} The whole connector sits inside the cell and the cable leaves through an
        epoxy-potted gland. Otherwise air diffuses along the jacket to the ferrule. Each cell must pass a
        pressure-decay leak test (target $<\SI{1e-5}{mbar.L/s}$, \cref{app:sops}); at that leak rate a
        typical absorber sachet lasts for months (\cref{app:estimates}).
  \item \textbf{Static variant (\arm{ABS}).} The absorber sits in a pocket behind a porous PTFE
        membrane, out of line of sight of the end face, so that no dust or vapour from the sachet
        reaches it. The cell is first purged with \Ntwo{} to about \SI{1}{\percent} \Oxy{}. The absorber
        then only has to remove a small residue, and the pressure in the sealed cell drops by about
        \SI{10}{mbar} instead of \SI{210}{mbar}, which would otherwise draw air in through any leak.
  \item \textbf{Verification.} An optical \Oxy{} sensor spot (read through a window, so that no
        electronics sit inside the cell) and a colour \Oxy{} indicator monitor the atmosphere.
  \item \textbf{Clean materials.} Printed parts are a possible source of organic vapour (H2). Prototypes
        are printed in PETG or ASA with \SI{100}{\percent} infill and several perimeters, sealed with a
        low-outgassing epoxy and baked under \Ntwo{} purge. The \arm{N2+P} arm tests directly whether the
        material harms the end face. O-rings are fitted dry or with a minimal film of PFPE grease, and
        silicone grease is banned because of siloxane contamination. The gland is potted with a
        low-outgassing epoxy.
\end{enumerate}
The student designs the cells in CAD with the dimensions of the group's fibre mounts. Two designs are
worth prototyping early: (a) a cap that encloses the connector and screws onto a standard
fibre-launch receptacle, and (b) a box that encloses the whole launch (lens included). Design (a) has a
smaller volume; design (b) is easier to seal.

\subsection{Diagnostics}
\begin{table}[t]
\centering\small
\caption{Measurements, their purpose and cadence.}
\label{tab:diagnostics}
\begin{tabularx}{\textwidth}{@{}L{4.1cm} Y L{2.6cm}@{}}
\toprule
Measurement & What it tells us & When \\
\midrule
$\eta = P_\text{out}/P_\text{ref}$ (normalised) & total throughput change of each arm & continuous, \SI{1}{Hz}\\
re-peak gain $\eta_\text{after}/\eta_\text{before}$ & alignment drift (H6) vs end-face change & daily or automated \\
$P_\text{back}/P_\text{ref}$ & optical change of the input face (film, roughness) & continuous \\
\Oxy, RH, VOC index, $T$, $p$, flow & that each atmosphere is what it should be & continuous \\
output beam profile and power through a mode-selective aperture & loss of beam quality (rises in $M^2$ have
been seen at \SI{405}{nm}~\cite{Gonschior2013PhD}) that a large detector does not register & every \SI{30}{min} \\
visible luminescence from the input face (spectrometer, optional) & fluorescence of growing organic films; defect PL bands & every \SI{30}{min} \\
end-face microscopy (bright field, fluorescence) & deposits, particles, damage & start and end of each run \\
cleaning sequence (\cref{sec:ladder}) & splits the loss between faces and bulk & end of each run \\
SEM/EDX, AFM or interferometry, Raman, ToF-SIMS/XPS & composition and thickness of films & selected ends \\
\bottomrule
\end{tabularx}
\end{table}
\Cref{tab:diagnostics} summarises the measurements. Two of them are cheap but very useful. First, the
\emph{power reflected by the input face} is \SI{3.7}{\percent} of the incident power. With the
\ang{8} FC/APC faces of the group's cables this reflection does not retrace the input beam. It leaves
\ang{\approx 24} away from it (\cref{app:estimates}) and can be caught with a small mirror or
photodiode. Even a thin film changes it through thin-film interference and scattering, so it is a
sensitive in-situ monitor of the input face. Second, \emph{organic films fluoresce} under near-UV excitation, while clean silica
hardly does. A long-pass filter and a small spectrometer that looks back through the coupling lens can
therefore see a film start to grow long before the throughput drops noticeably.

\subsubsection{End-of-run cleaning sequence}\label{sec:ladder}
At the end of each run the transmission is measured at low power after each of the following steps,
always with a re-peak and first at the output end, then at the input end:
\begin{enumerate}
  \item no-touch inspection and imaging (preserves evidence);
  \item dry clean (filtered gas jet, then a dry cleaning tool);
  \item solvent clean (reagent-grade isopropanol, lint-free wipe; acetone as a second solvent);
  \item oxidative clean (\Oxy{} plasma or UV/ozone, if available), which removes organic but not
        inorganic (e.g.\ siloxane-derived \ce{SiO_x}) films;
  \item re-polish, recording the removed length.
\end{enumerate}
The loss recovered at each step assigns the damage to a face and to a category (loose particles,
soluble film, oxidisable film, bonded or structural damage). Any loss left after both ends have been
re-polished is attributed to the fibre itself (H5). Where an arm has two fibres, one goes through
this sequence and the other is kept uncleaned for surface analysis.

\subsection{Experimental campaigns}\label{sec:campaigns}
\begin{description}[style=nextline,leftmargin=1.2em,itemsep=4pt]
\item[Phase 1 -- baseline in ambient air (months 4--5)] Three to four \arm{AIR} arms at \SI{1}{mW},
  \SI{3}{mW} and \SI{10}{mW}, plus the sentinel. Outputs: (i) proof that the testbed reproduces the
  degradation seen in the lab; (ii) the power law $r\propto P^{n}$ of the initial degradation rate. An
  exponent $n\approx2$ points to a two-photon process; (iii) fibre-to-fibre scatter; (iv) the run length
  needed in the later campaigns. Raising the power only accelerates the test if the mechanism stays the
  same. Surface restructuring at \SI{405}{nm} was reported above \SI{\sim0.5}{MW/cm^2}~\cite{Gonschior2013JNCS},
  close to the \SI{0.44}{MW/cm^2} reached here at \SI{10}{mW}. The end faces from the \SI{3}{mW} and
  \SI{10}{mW} arms are therefore compared by microscopy before \SI{10}{mW} is used as an accelerated condition.
\item[Phase 2 -- Campaign I, atmosphere screening (months 6--7)] \arm{AIR}, \arm{CDA}, \arm{N2} and
  \arm{ABS} in parallel at the accelerated power (or \SI{3}{mW} if Phase~1 shows that \SI{10}{mW} changes
  the mechanism, e.g.\ thermal runaway). Each run lasts until \arm{AIR} has lost \SI{\sim30}{\percent} or
  for three weeks, whichever comes first. If time allows, a second run with permuted arm positions follows.
\item[Phase 3 -- Campaign II, confirmation and mechanism (months 8--9)] Repeat the key contrast at
  \SI{3}{mW} with at least two fibres per condition, and add the most informative extended arms
  (\cref{tab:arms}): the \Oxy{} dose--response if \Oxy{} matters, \arm{WET-N2}/\arm{WET-CDA} and \arm{N2+P}
  if it does not. A crossed configuration (input in \Ntwo{}, output in air, and vice versa) tests whether
  both faces respond in the same way. Marciniak et al.\ observed that the purged and unpurged ends of
  the same fibre behaved differently~\cite{Marciniak2017}.
\item[Phase 4 -- demonstration at 369.5\,nm (months 9--11)] The best cell design is fitted to both ends
  of a working \SI{369.5}{nm} fibre in one of the group's experiments, with a control fibre or the fibre's
  own history as reference, for at least eight weeks. This phase also checks the transfer from
  \SI{375}{nm} to \SI{369.5}{nm} and whether the cell is practical in use (alignment stability, access,
  window losses).
\end{description}

\begin{decisionbox}[How the outcome of Campaign I steers the project]
\begin{itemize}[leftmargin=1em]
  \item \textbf{\arm{N2} and \arm{ABS} stable, \arm{CDA} degrades like \arm{AIR}}: oxygen takes part
        (H1). Check the deposit chemistry: if it is rich in Si and O, siloxanes are involved, and their
        sources near the fibres should be removed anyway. Then measure the \Oxy{} dose--response,
        specify the cell atmosphere and deploy \arm{ABS} cells. This is the outcome the group hopes for.
  \item \textbf{\arm{CDA} and \arm{N2} both stable, \arm{AIR} degrades}: something in lab air other than
        \Oxy{} is responsible: particles, organic vapours or water (H4, H2, or a water-assisted H3).
        Campaign~II separates them. A sealed or purged \emph{clean} cell solves the practical problem, with
        or without oxygen. This is what the published nitrogen-purge results would suggest~\cite{Marciniak2017}.
  \item \textbf{\arm{CDA} stable, \arm{N2} or \arm{ABS} degrade}: carbonaceous contamination with
        oxygen as a protective agent (H2). This is the outcome that space-laser, UV-cavity and laser-diode
        experience predicts (\cref{sec:sota-surface}). Oxygen absorbers would be counter-productive; the
        recommendation becomes a clean, dry (or humidified) air purge.
  \item \textbf{No atmosphere helps}: restructuring of the silica, colour centres or an artefact (H3, H5,
        H6). Shift the effort to end-face and fibre engineering. Coreless end caps lower the intensity at
        the glass--gas interface by two to three orders of magnitude (\cref{app:estimates}). Other options are
        hydrogen-loaded and UV-cured fibre and hollow-core fibre (\cref{sec:sota-fibres}).
\end{itemize}
Every branch gives a publishable answer and a concrete recommendation for the group.
\end{decisionbox}

\subsection{Data analysis and decision rules}\label{sec:analysis}
\paragraph{Metrics} The primary quantity is the normalised, re-peaked throughput
$\eta(t)=[P_\text{out}/P_\text{ref}](t)/[P_\text{out}/P_\text{ref}](t_0)$, corrected for detector
temperature coefficients, with flagged intervals (interventions, outages) excluded. The \emph{primary
endpoint} is the initial degradation rate $r=-\mathrm{d}\eta/\mathrm{d}t$ in \si{\percent} per
\SI{1000}{h} at \SI{3}{mW}. Secondary endpoints are the time to \SI{10}{\percent} loss, the fit
parameters of a phenomenological model (e.g.\ a stretched exponential
$\eta=1-A\{1-\exp[-(t/\tau)^{\beta}]\}$ in time or dose), the change of $P_\text{back}$, and the outcome of
the cleaning sequence.

\paragraph{Statistics} Confidence intervals come from a block bootstrap, which allows for
autocorrelated data. Conditions are compared through rate ratios. With replicate fibres, a
mixed-effects model with fibre as random effect separates the atmosphere effect from fibre-to-fibre
scatter. Phase~1 provides the scatter needed to choose the number of replicates.

\paragraph{Decision rules (fixed before Campaign II)}
\begin{itemize}
  \item \textbf{\Oxy{} effect established:} $r_\text{N2}/r_\text{CDA}$ differs from 1 by at least a factor of
        two, with the 95\,\% confidence interval excluding 1, in two independent runs.
  \item \textbf{Degradation halted:} the 95\,\% upper limit on $|r|$ is below \SI{1}{\percent} per
        \SI{1000}{h} at \SI{3}{mW} (equivalent to $<\SI{10}{\percent}$ loss per year of continuous use),
        at least twenty times below \arm{AIR}, with no visible change of either end face. With a daily
        scatter of \SI{0.3}{\percent} this limit can be reached in three to six weeks (\cref{app:estimates}).
\end{itemize}

\subsection{Success criteria and deliverables}
\begin{enumerate}[label=\textbf{S\arabic*}]
  \item The location and physical nature of the degradation are established for the group's fibres (RQ1).
  \item The effect of \Oxy{} on the degradation rate is measured with a stated confidence interval, or
        excluded at the level above (RQ2), and the roles of water, contaminants and particles are tested (RQ3).
  \item A cell design is qualified: \Oxy{} $<\SI{100}{ppm}$ (target $<\SI{10}{ppm}$) and RH
        $<\SI{5}{\percent}$ for at least eight weeks without maintenance. It is tested on a working
        \SI{369.5}{nm} link (RQ4). If the atmosphere helps, the degradation-halted criterion is met.
  \item Open deliverables: CAD files and bill of materials of the cells, standard operating procedures, the
        data set with analysis code, the MSc thesis, and (if the results allow) a short paper.
\end{enumerate}
